\section{The Entity Hop and its indexing cost}
\label{sec:hop}

\paragraph{Indexing unit and first stage.} The unit is one whole article paragraph, with its title
prefixed; never a fixed-size window. Dense is BGE-small~\citep{xiao2024cpack} over
L2-normalized embeddings, and BM25~\citep{robertson2009bm25} is the lexical complement. Two
hybrids are fused by min-max normalizing each component's scores per question and taking a
weighted sum, with the weights fitted on HotpotQA development questions and then frozen.

\paragraph{The hop.} Let $D(q)$ be the Dense ranking for question $q$. We read the first
ten units of $D(q)$ and call the first unit $p_1(q)$. The candidates are the paragraphs
\emph{outside} the units read that share at least one raw entity node with $p_1(q)$. The score
of a candidate is the sum, over the entity nodes it shares with $p_1(q)$, of the rarity
\begin{equation}
  \log\!\Bigl(1+\tfrac{N}{1+\mathrm{df}}\Bigr),
\end{equation}
where $N$ is the number of units and $\mathrm{df}$ the number of units holding the entity.
Candidates are ranked by score, then by unit id; the list is not padded. The hop never looks at
the question after the first Dense call, uses one seed and one hop, and has no relations, no
canonicalization and no language model at query time. Entities are read once, at indexing,
by GLiNER~\citep{zaratiana2024gliner} (the medium model, zero-shot, five entity
labels).

\paragraph{Four systems.} Throughout, \textbf{P10-A} is Dense, \textbf{P10-B} is Dense plus BM25,
\textbf{P10-C} is Dense plus BM25 plus the Entity Hop as a third fused component, and
\textbf{P14} is P10-C with a single change: the hop scores its candidates by
$(1-\alpha)\cdot\text{rarity}+\alpha\cdot\cos(q,c)$, both terms min-max normalized per question,
with $\alpha=\phFourteenAlpha$ fitted on development questions. The seed, the candidate set and
the fusion weights of P14 are P10-C's; only the order of the hop's candidates changes. The
similarity is Dense's own signal, so part of P14's gain may re-weight what Dense already ranks;
on development data, among the gold paragraphs whose coverage changed in the questions P14 won
against P10-C, about half were in Dense's ranks 11 to 100 and about half were beyond its first
hundred (derived, Phase~14); the split cannot separate the two causes.

\paragraph{What did not work as a refinement.} Three attempts to make the hop choose which of
$p_1$'s entities to follow, with a hard filter (a document-frequency cap, entities named in the
question, a window around each mention), each lost to P10-C on development questions and stopped
at their gate, so the held-out split was never opened for them. Phase~14 then kept all of the
hop's reach and changed only its order, and that is the version that held. The reading we take
from this is a hypothesis, not an established cause: the hard filter, not the signal, is what
fails.

\begin{table}[t]
\centering
\small
\caption{\capHopCost}
\label{tab:hopcost}
\coltable{tables/hop-cost}
\end{table}

\paragraph{Indexing cost.} Table~\ref{tab:hopcost} gives the cost of one entity-extraction
pass. Over the whole HotpotQA Wikipedia (\hopCostNineUnits{} paragraphs) GLiNER took
\hopCostNineHours{} hours on one rented RTX~4090, for an attributable \hopCostNineUsd{} USD
(derived from the measured time and the hourly rate), plus \hopCostNineEmbedUsd{} USD to embed
the same paragraphs. The same extractor over the MuSiQue and MultiHop-RAG corpora took
\hopCostFifteenUsd{} and \hopCostSixteenUsd{} USD. On the small pool where both were run,
the language-model extractor (the Claude row) cost \phSevenClaudeUsd{} USD against \phSevenGlinerUsd{}
USD for GLiNER, a ratio of about \phSevenCostRatio{} (derived), and GLiNER kept \phSevenRetention{}
of the language-model extractor's gain. Extracting the full corpus with the language model would
cost on the order of \phSevenClaudeProjectedUsd{} USD; that is a \emph{projection} from linear
scaling, not a measurement. At query time the hop adds a lookup and a fusion; on the laptop,
P10-C's mean latency was \phTenLatencyRatio{} times that of Dense plus BM25 (measured, Phase~10).
